\chapter*{新版前言：从智能结构到状态动力学}
\addcontentsline{toc}{chapter}{新版前言：从智能结构到状态动力学}
\markboth{前言}{从智能结构到状态动力学}
本书研究一个贯穿智能理论与工程的问题：一个具有有限资源的系统，如何在持续交互中组织当前活动、积累历史、改变结构，并维持目标导向的行动？我们将这一问题表述为智能状态的计算动力学。

我们保留全息语义场与层级动力学（HSF-HD）的名称。“语义场”表示在关系结构上分布并相互作用的语义状态；“全息”表示局部活动受到整体上下文约束，并能够参与重建任务相关的整体信息；“层级动力学”表示不同作用范围与更新速度的过程相互耦合。这些术语是计算建模的约定，不预设认知过程遵循某种量子或引力方程。

本书的理论位置介于基础系统科学与智能运行时架构之间。系统论、协同论、耗散结构理论和复杂系统理论提供组织、反馈、开放性与多尺度演化的视角，认知科学提供记忆、注意、控制和自我模型的经验约束。HSF-HD 将这些约束组织为可以分析和验证的模型族，AgentOS 则将模型实例化为记忆管理、状态监控、资源调度、工具交互与持续学习机制。

我们把“智能的一般空气动力学”作为研究方向的类比：如同研究飞行需要区分运动状态、结构、边界和控制，研究智能也需要区分当前活动、持久结构、输入输出、控制策略与资源预算。类比帮助提出问题，结论仍须由明确假设下的数学推导和实现中的观测支持。

全书保留三卷十一篇。第一卷建立结构与属性的表示，第二卷研究状态、记忆、控制和多尺度耦合，第三卷比较不同系统并讨论工程实现。几何方法是可选的表示工具，物理学语言用于说明载体约束或提供辅助直觉。计算量、信息量、状态范数和实际能耗分别定义，不能相互替代。

我们使用“定义”建立模型语言，使用“假设”声明适用条件，使用“命题”给出能够推导的结果。尚缺验证的机制保持为模型假说。一个方程可以在数学上成立而在经验上不适用；一个原型可以执行而尚未具有理论预期的性质。这些区分构成理论走向工程的必要步骤。

本书的目标不是用一种术语覆盖所有智能现象，而是建立一套可以逐步修正的共同语言，使状态演化、结构学习和目标控制能够在同一框架中被讨论、实现与检验。
\hfill XuEn

我们保留序言中的文学表达和原有章节的问题展开，以延续本书的思想来源。后续论证以计算动力学中的对象类型、作用机制与可检验条件为准。将模型用于某一认知现象时，我们需要另行给出观测、比较对象和证据，不能把数学表达本身视为现象已经得到解释。
